% --- Archon physics formalization source begin ---
% archon:physics
% archon:covers PhyXMiniProblems/problem_phyx_mini_0995.lean
% archon:source-report reports/phyx_mini/problem_phyx_mini_0995.source.json

\chapter{Physics problem phyx\_mini\_0995}
\label{ch:PhyXMiniProblems_problem_phyx_mini_0995}

\paragraph{Problem source.}
\#\# Physical scenario

Charge \textbackslash{}( Q \textbackslash{}) is uniformly distributed around a conducting ring of radius \textbackslash{}( a \textbackslash{}).

\#\# Question

Find the electric field at a point \textbackslash{}( P \textbackslash{}) on the ring axis at a distance \textbackslash{}( x \textbackslash{}) from its center.

\#\# Answer choices

- A: A: \textbackslash{}frac\{1\}\{4\textbackslash{}pi\textbackslash{}varepsilon\_0\} \textbackslash{}frac\{Qa\}\{(x\textasciicircum{}2 + a\textasciicircum{}2)\textasciicircum{}\{3/2\}\} \textbackslash{}hat\{i\}

- B: B: \textbackslash{}frac\{1\}\{4\textbackslash{}pi\textbackslash{}varepsilon\_0\} \textbackslash{}frac\{Qx\}\{(x\textasciicircum{}2 + a\textasciicircum{}2)\textasciicircum{}\{3/2\}\} \textbackslash{}hat\{i\}

- C: C: \textbackslash{}frac\{1\}\{4\textbackslash{}pi\textbackslash{}varepsilon\_0\} \textbackslash{}frac\{Qx\}\{(x\textasciicircum{}2 + a\textasciicircum{}2)\textasciicircum{}\{1/2\}\} \textbackslash{}hat\{i\}

- D: D: \textbackslash{}frac\{1\}\{4\textbackslash{}pi\textbackslash{}varepsilon\_0\} \textbackslash{}frac\{Q\}\{(x\textasciicircum{}2 + a\textasciicircum{}2)\} \textbackslash{}hat\{i\}

\#\# Figure caption (auxiliary; use the image as primary evidence)

The image depicts a diagram of a charged ring in a coordinate system. Here's a detailed description:

- **Coordinate System**: The diagram includes x and y axes intersecting at point \textbackslash{}(O\textbackslash{}).
- **Charged Ring**: A ring centered at \textbackslash{}(O\textbackslash{}) with radius \textbackslash{}(a\textbackslash{}) is shown. It is labeled with a total charge \textbackslash{}(Q\textbackslash{}).
- **Element of Charge**: A small segment \textbackslash{}(ds\textbackslash{}) on the ring is highlighted, associated with an infinitesimal charge \textbackslash{}(dQ\textbackslash{}).
- **Point \textbackslash{}(P\textbackslash{})**: Located on the x-axis at a distance from the origin. This is where the electric field is being evaluated.
- **Vectors and Angles**: 
  - A dashed line \textbackslash{}(r\textbackslash{}) connects the charge element \textbackslash{}(dQ\textbackslash{}) to point \textbackslash{}(P\textbackslash{}), with \textbackslash{}(r = \textbackslash{}sqrt\{x\textasciicircum{}2 + a\textasciicircum{}2\}\textbackslash{}).
  - The angle between the x-axis and line \textbackslash{}(r\textbackslash{}) is \textbackslash{}(\textbackslash{}alpha\textbackslash{}).
  - The electric field components due to \textbackslash{}(dQ\textbackslash{}) at \textbackslash{}(P\textbackslash{}) are shown as \textbackslash{}(dE\_x\textbackslash{}) (horizontal) and \textbackslash{}(dE\_y\textbackslash{}) (vertical) with \textbackslash{}(dE\textbackslash{}) as the resultant electric field vector. 
  - The angle \textbackslash{}(\textbackslash{}theta\textbackslash{}) is given between \textbackslash{}(dE\textbackslash{}) and \textbackslash{}(dE\_x\textbackslash{}).

The diagram illustrates the contributions to the electric field from an element \textbackslash{}(dQ\textbackslash{}) of the charged ring at a given point along the x-axis.

\#\# Dataset metadata

Electrostatics; Physical Model Grounding Reasoning, Spatial Relation Reasoning

\paragraph{Recorded answer/context.}
B: B: \textbackslash{}frac\{1\}\{4\textbackslash{}pi\textbackslash{}varepsilon\_0\} \textbackslash{}frac\{Qx\}\{(x\textasciicircum{}2 + a\textasciicircum{}2)\textasciicircum{}\{3/2\}\} \textbackslash{}hat\{i\}

\paragraph{Figure/image path.}
/root/proposal\_for\_physic/science-mango/phyx\_mini\_run/phyx\_data/test\_image/995.png

\paragraph{Formalization target.}
create a compiling Lean file with sorry bodies at `PhyXMiniProblems/problem\_phyx\_mini\_0995.lean`. The Lean declarations must preserve the physical quantities, dimensions or dimensional roles, figure labels, governing-law hypotheses, and final relation expressed by this problem.
Use Mathlib/Physlib names found through LeanExplore where available. If a physics API is missing, introduce faithful local abstractions rather than scalar placeholder aliases.

\begin{theorem}[Physics formalization target]
\label{thm:physics:phyx_mini_0995:target}
\lean{PhyXMiniProblems.ProblemPhyXMini0995.electricFieldOnAxisOfUniformlyChargedRing}
\uses{def:physics:phyx-mini-0995:phyxminiproblems-problemphyxmini0995-chargeincoulombs, def:physics:phyx-mini-0995:phyxminiproblems-problemphyxmini0995-lengthinmeters, def:physics:phyx-mini-0995:phyxminiproblems-problemphyxmini0995-chargedringaxissetup, def:physics:phyx-mini-0995:phyxminiproblems-problemphyxmini0995-xaxisunit, def:physics:phyx-mini-0995:phyxminiproblems-problemphyxmini0995-matchesconductingringscenario, def:physics:phyx-mini-0995:phyxminiproblems-problemphyxmini0995-matchesprimarychargedringfigure, def:physics:phyx-mini-0995:phyxminiproblems-problemphyxmini0995-haschargedringaxisgeometry, def:physics:phyx-mini-0995:phyxminiproblems-problemphyxmini0995-hasphysicalelectrostaticparameters, def:physics:phyx-mini-0995:phyxminiproblems-problemphyxmini0995-satisfiesuniformringdistribution, def:physics:phyx-mini-0995:phyxminiproblems-problemphyxmini0995-satisfiesdifferentialcoulombfieldlaw, def:physics:phyx-mini-0995:phyxminiproblems-problemphyxmini0995-satisfiesringfieldsuperposition}
The assigned autoformalize agent should translate this physics problem into Lean declarations in the covered file, with theorem and lemma proof bodies written as `by sorry`.
\end{theorem}
\begin{proof}
This is an autoformalization task, not a proof task. Produce statements that can later be proved by the physics prover without weakening the physical model.
\end{proof}

% --- Archon declaration topology begin ---
\section{Declaration topology}
\begin{definition}[Signed Charge Quantity]
  \label{def:physics:phyx-mini-0995:phyxminiproblems-problemphyxmini0995-signedchargequantity}
  \lean{PhyXMiniProblems.ProblemPhyXMini0995.SignedChargeQuantity}
  A signed, unit-independent physical electric charge
\end{definition}
\begin{proof}
  This is the declared model definition of Signed Charge Quantity.
\end{proof}
\begin{definition}[Length Quantity]
  \label{def:physics:phyx-mini-0995:phyxminiproblems-problemphyxmini0995-lengthquantity}
  \lean{PhyXMiniProblems.ProblemPhyXMini0995.LengthQuantity}
  A nonnegative, unit-independent physical length
\end{definition}
\begin{proof}
  This is the declared model definition of Length Quantity.
\end{proof}
\begin{definition}[Angular Charge Density Quantity]
  \label{def:physics:phyx-mini-0995:phyxminiproblems-problemphyxmini0995-angularchargedensityquantity}
  \lean{PhyXMiniProblems.ProblemPhyXMini0995.AngularChargeDensityQuantity}
  An angular charge density has charge dimension because radians are dimensionless. Its SI readout has the role coulombs per radian
\end{definition}
\begin{proof}
  This is the declared model definition of Angular Charge Density Quantity.
\end{proof}
\begin{definition}[charge In Coulombs]
  \label{def:physics:phyx-mini-0995:phyxminiproblems-problemphyxmini0995-chargeincoulombs}
  \lean{PhyXMiniProblems.ProblemPhyXMini0995.chargeInCoulombs}
  \uses{def:physics:phyx-mini-0995:phyxminiproblems-problemphyxmini0995-signedchargequantity}
  Read a signed charge in coherent-SI coulombs
\end{definition}
\begin{proof}
  This is the declared model definition of charge In Coulombs.
\end{proof}
\begin{definition}[length In Meters]
  \label{def:physics:phyx-mini-0995:phyxminiproblems-problemphyxmini0995-lengthinmeters}
  \lean{PhyXMiniProblems.ProblemPhyXMini0995.lengthInMeters}
  \uses{def:physics:phyx-mini-0995:phyxminiproblems-problemphyxmini0995-lengthquantity}
  Read a physical length in coherent-SI metres
\end{definition}
\begin{proof}
  This is the declared model definition of length In Meters.
\end{proof}
\begin{definition}[angular Charge Density In Coulombs Per Radian]
  \label{def:physics:phyx-mini-0995:phyxminiproblems-problemphyxmini0995-angularchargedensityincoulombsperradian}
  \lean{PhyXMiniProblems.ProblemPhyXMini0995.angularChargeDensityInCoulombsPerRadian}
  \uses{def:physics:phyx-mini-0995:phyxminiproblems-problemphyxmini0995-angularchargedensityquantity}
  Read angular charge density in coherent-SI coulombs per radian
\end{definition}
\begin{proof}
  This is the declared model definition of angular Charge Density In Coulombs Per Radian.
\end{proof}
\begin{definition}[Figure Axis]
  \label{def:physics:phyx-mini-0995:phyxminiproblems-problemphyxmini0995-figureaxis}
  \lean{PhyXMiniProblems.ProblemPhyXMini0995.FigureAxis}
  Coordinate axes explicitly drawn in the supplied raster
\end{definition}
\begin{proof}
  This is the declared model definition of Figure Axis.
\end{proof}
\begin{definition}[Figure Point]
  \label{def:physics:phyx-mini-0995:phyxminiproblems-problemphyxmini0995-figurepoint}
  \lean{PhyXMiniProblems.ProblemPhyXMini0995.FigurePoint}
  Named points or highlighted source elements in the supplied raster
\end{definition}
\begin{proof}
  This is the declared model definition of Figure Point.
\end{proof}
\begin{definition}[Figure Scalar Label]
  \label{def:physics:phyx-mini-0995:phyxminiproblems-problemphyxmini0995-figurescalarlabel}
  \lean{PhyXMiniProblems.ProblemPhyXMini0995.FigureScalarLabel}
  Scalar or differential labels printed in the supplied raster
\end{definition}
\begin{proof}
  This is the declared model definition of Figure Scalar Label.
\end{proof}
\begin{definition}[Figure Vector Label]
  \label{def:physics:phyx-mini-0995:phyxminiproblems-problemphyxmini0995-figurevectorlabel}
  \lean{PhyXMiniProblems.ProblemPhyXMini0995.FigureVectorLabel}
  Vector and vector-component labels printed in the supplied raster
\end{definition}
\begin{proof}
  This is the declared model definition of Figure Vector Label.
\end{proof}
\begin{definition}[Figure Angle Label]
  \label{def:physics:phyx-mini-0995:phyxminiproblems-problemphyxmini0995-figureanglelabel}
  \lean{PhyXMiniProblems.ProblemPhyXMini0995.FigureAngleLabel}
  Angle labels printed near the source line and field-component triangle
\end{definition}
\begin{proof}
  This is the declared model definition of Figure Angle Label.
\end{proof}
\begin{definition}[Figure Feature]
  \label{def:physics:phyx-mini-0995:phyxminiproblems-problemphyxmini0995-figurefeature}
  \lean{PhyXMiniProblems.ProblemPhyXMini0995.FigureFeature}
  Other geometric features visible in image 995.png
\end{definition}
\begin{proof}
  This is the declared model definition of Figure Feature.
\end{proof}
\begin{definition}[expected Axis Label]
  \label{def:physics:phyx-mini-0995:phyxminiproblems-problemphyxmini0995-expectedaxislabel}
  \lean{PhyXMiniProblems.ProblemPhyXMini0995.expectedAxisLabel}
  \uses{def:physics:phyx-mini-0995:phyxminiproblems-problemphyxmini0995-figureaxis}
  Literal text expected beside each drawn coordinate axis
\end{definition}
\begin{proof}
  This is the declared model definition of expected Axis Label.
\end{proof}
\begin{definition}[expected Point Label]
  \label{def:physics:phyx-mini-0995:phyxminiproblems-problemphyxmini0995-expectedpointlabel}
  \lean{PhyXMiniProblems.ProblemPhyXMini0995.expectedPointLabel}
  \uses{def:physics:phyx-mini-0995:phyxminiproblems-problemphyxmini0995-figurepoint}
  Literal text expected beside each named point or element
\end{definition}
\begin{proof}
  This is the declared model definition of expected Point Label.
\end{proof}
\begin{definition}[expected Scalar Label]
  \label{def:physics:phyx-mini-0995:phyxminiproblems-problemphyxmini0995-expectedscalarlabel}
  \lean{PhyXMiniProblems.ProblemPhyXMini0995.expectedScalarLabel}
  \uses{def:physics:phyx-mini-0995:phyxminiproblems-problemphyxmini0995-figurescalarlabel}
  Literal scalar/differential text transcribed from the image
\end{definition}
\begin{proof}
  This is the declared model definition of expected Scalar Label.
\end{proof}
\begin{definition}[expected Vector Label]
  \label{def:physics:phyx-mini-0995:phyxminiproblems-problemphyxmini0995-expectedvectorlabel}
  \lean{PhyXMiniProblems.ProblemPhyXMini0995.expectedVectorLabel}
  \uses{def:physics:phyx-mini-0995:phyxminiproblems-problemphyxmini0995-figurevectorlabel}
  Literal vector text transcribed from the image
\end{definition}
\begin{proof}
  This is the declared model definition of expected Vector Label.
\end{proof}
\begin{definition}[expected Angle Label]
  \label{def:physics:phyx-mini-0995:phyxminiproblems-problemphyxmini0995-expectedanglelabel}
  \lean{PhyXMiniProblems.ProblemPhyXMini0995.expectedAngleLabel}
  \uses{def:physics:phyx-mini-0995:phyxminiproblems-problemphyxmini0995-figureanglelabel}
  Literal Greek angle text transcribed from the image
\end{definition}
\begin{proof}
  This is the declared model definition of expected Angle Label.
\end{proof}
\begin{definition}[Charged Ring Axis Figure]
  \label{def:physics:phyx-mini-0995:phyxminiproblems-problemphyxmini0995-chargedringaxisfigure}
  \lean{PhyXMiniProblems.ProblemPhyXMini0995.ChargedRingAxisFigure}
  \uses{def:physics:phyx-mini-0995:phyxminiproblems-problemphyxmini0995-signedchargequantity, def:physics:phyx-mini-0995:phyxminiproblems-problemphyxmini0995-lengthquantity, def:physics:phyx-mini-0995:phyxminiproblems-problemphyxmini0995-figureaxis, def:physics:phyx-mini-0995:phyxminiproblems-problemphyxmini0995-figurepoint, def:physics:phyx-mini-0995:phyxminiproblems-problemphyxmini0995-figurescalarlabel, def:physics:phyx-mini-0995:phyxminiproblems-problemphyxmini0995-figurevectorlabel, def:physics:phyx-mini-0995:phyxminiproblems-problemphyxmini0995-figureanglelabel, def:physics:phyx-mini-0995:phyxminiproblems-problemphyxmini0995-figurefeature}
  Presentation data and physical label calibrations read from the primary image. The angle values describe the one highlighted element only; no unwarranted equality between the visibly distinct α and θ marks is made
\end{definition}
\begin{proof}
  This is the declared model definition of Charged Ring Axis Figure.
\end{proof}
\begin{definition}[Charged Ring Source]
  \label{def:physics:phyx-mini-0995:phyxminiproblems-problemphyxmini0995-chargedringsource}
  \lean{PhyXMiniProblems.ProblemPhyXMini0995.ChargedRingSource}
  \uses{def:physics:phyx-mini-0995:phyxminiproblems-problemphyxmini0995-signedchargequantity, def:physics:phyx-mini-0995:phyxminiproblems-problemphyxmini0995-lengthquantity, def:physics:phyx-mini-0995:phyxminiproblems-problemphyxmini0995-angularchargedensityquantity}
  The source curve is parametrized by a dimensionless angle in radians. The two density fields represent the figure's differential elements through dQ = (dQ/dφ) dφ and ds = (ds/dφ) dφ
\end{definition}
\begin{proof}
  This is the declared model definition of Charged Ring Source.
\end{proof}
\begin{definition}[Charged Ring Axis Setup]
  \label{def:physics:phyx-mini-0995:phyxminiproblems-problemphyxmini0995-chargedringaxissetup}
  \lean{PhyXMiniProblems.ProblemPhyXMini0995.ChargedRingAxisSetup}
  \uses{def:physics:phyx-mini-0995:phyxminiproblems-problemphyxmini0995-lengthquantity, def:physics:phyx-mini-0995:phyxminiproblems-problemphyxmini0995-chargedringaxisfigure, def:physics:phyx-mini-0995:phyxminiproblems-problemphyxmini0995-chargedringsource}
  The electric field and its differential angular contributions are primitive observables. Their relation is imposed later by Coulomb and superposition laws; neither is defined from the requested answer
\end{definition}
\begin{proof}
  This is the declared model definition of Charged Ring Axis Setup.
\end{proof}
\begin{definition}[x Axis Unit]
  \label{def:physics:phyx-mini-0995:phyxminiproblems-problemphyxmini0995-xaxisunit}
  \lean{PhyXMiniProblems.ProblemPhyXMini0995.xAxisUnit}
  Unit vector î along the displayed positive x-axis
\end{definition}
\begin{proof}
  This is the declared model definition of x Axis Unit.
\end{proof}
\begin{definition}[source To Observation Displacement In Meters]
  \label{def:physics:phyx-mini-0995:phyxminiproblems-problemphyxmini0995-sourcetoobservationdisplacementinmeters}
  \lean{PhyXMiniProblems.ProblemPhyXMini0995.sourceToObservationDisplacementInMeters}
  \uses{def:physics:phyx-mini-0995:phyxminiproblems-problemphyxmini0995-chargedringaxissetup}
  Coherent-SI displacement from a ring source element to P
\end{definition}
\begin{proof}
  This is the declared model definition of source To Observation Displacement In Meters.
\end{proof}
\begin{definition}[source To Observation Distance In Meters]
  \label{def:physics:phyx-mini-0995:phyxminiproblems-problemphyxmini0995-sourcetoobservationdistanceinmeters}
  \lean{PhyXMiniProblems.ProblemPhyXMini0995.sourceToObservationDistanceInMeters}
  \uses{def:physics:phyx-mini-0995:phyxminiproblems-problemphyxmini0995-chargedringaxissetup, def:physics:phyx-mini-0995:phyxminiproblems-problemphyxmini0995-sourcetoobservationdisplacementinmeters}
  Coherent-SI source-to-P distance labelled r in the figure
\end{definition}
\begin{proof}
  This is the declared model definition of source To Observation Distance In Meters.
\end{proof}
\begin{definition}[Matches Conducting Ring Scenario]
  \label{def:physics:phyx-mini-0995:phyxminiproblems-problemphyxmini0995-matchesconductingringscenario}
  \lean{PhyXMiniProblems.ProblemPhyXMini0995.MatchesConductingRingScenario}
  \uses{def:physics:phyx-mini-0995:phyxminiproblems-problemphyxmini0995-chargedringaxissetup}
  The physical source has the conducting-ring role stated in the scenario
\end{definition}
\begin{proof}
  This is the declared model definition of Matches Conducting Ring Scenario.
\end{proof}
\begin{definition}[Matches Primary Charged Ring Figure]
  \label{def:physics:phyx-mini-0995:phyxminiproblems-problemphyxmini0995-matchesprimarychargedringfigure}
  \lean{PhyXMiniProblems.ProblemPhyXMini0995.MatchesPrimaryChargedRingFigure}
  \uses{def:physics:phyx-mini-0995:phyxminiproblems-problemphyxmini0995-expectedaxislabel, def:physics:phyx-mini-0995:phyxminiproblems-problemphyxmini0995-expectedpointlabel, def:physics:phyx-mini-0995:phyxminiproblems-problemphyxmini0995-expectedscalarlabel, def:physics:phyx-mini-0995:phyxminiproblems-problemphyxmini0995-expectedvectorlabel, def:physics:phyx-mini-0995:phyxminiproblems-problemphyxmini0995-expectedanglelabel, def:physics:phyx-mini-0995:phyxminiproblems-problemphyxmini0995-chargedringaxissetup}
  Literal labels and qualitative features of the primary raster, calibrated to the independent physical quantities. The source-distance formula is retained as displayed text here and derived mathematically for every source angle below
\end{definition}
\begin{proof}
  This is the declared model definition of Matches Primary Charged Ring Figure.
\end{proof}
\begin{definition}[Has Charged Ring Axis Geometry]
  \label{def:physics:phyx-mini-0995:phyxminiproblems-problemphyxmini0995-haschargedringaxisgeometry}
  \lean{PhyXMiniProblems.ProblemPhyXMini0995.HasChargedRingAxisGeometry}
  \uses{def:physics:phyx-mini-0995:phyxminiproblems-problemphyxmini0995-lengthinmeters, def:physics:phyx-mini-0995:phyxminiproblems-problemphyxmini0995-chargedringaxissetup}
  The displayed x-axis is coordinate 0; the ring lies in the orthogonal y-z plane and is parametrized by (a cos φ, a sin φ). Thus O is the ring center and P is x metres along the positive axis
\end{definition}
\begin{proof}
  This is the declared model definition of Has Charged Ring Axis Geometry.
\end{proof}
\begin{definition}[Has Physical Electrostatic Parameters]
  \label{def:physics:phyx-mini-0995:phyxminiproblems-problemphyxmini0995-hasphysicalelectrostaticparameters}
  \lean{PhyXMiniProblems.ProblemPhyXMini0995.HasPhysicalElectrostaticParameters}
  \uses{def:physics:phyx-mini-0995:phyxminiproblems-problemphyxmini0995-chargeincoulombs, def:physics:phyx-mini-0995:phyxminiproblems-problemphyxmini0995-chargedringaxissetup}
  Nonzero source charge and a physically positive Coulomb constant
\end{definition}
\begin{proof}
  This is the declared model definition of Has Physical Electrostatic Parameters.
\end{proof}
\begin{definition}[Satisfies Uniform Ring Distribution]
  \label{def:physics:phyx-mini-0995:phyxminiproblems-problemphyxmini0995-satisfiesuniformringdistribution}
  \lean{PhyXMiniProblems.ProblemPhyXMini0995.SatisfiesUniformRingDistribution}
  \uses{def:physics:phyx-mini-0995:phyxminiproblems-problemphyxmini0995-chargeincoulombs, def:physics:phyx-mini-0995:phyxminiproblems-problemphyxmini0995-lengthinmeters, def:physics:phyx-mini-0995:phyxminiproblems-problemphyxmini0995-angularchargedensityincoulombsperradian, def:physics:phyx-mini-0995:phyxminiproblems-problemphyxmini0995-chargedringaxissetup}
  Uniform distribution means dQ/dφ = Q/(2π). The corresponding arc-length density is ds/dφ = a; the integral clauses explicitly account for total charge and circumference without assuming the requested electric field
\end{definition}
\begin{proof}
  This is the declared model definition of Satisfies Uniform Ring Distribution.
\end{proof}
\begin{definition}[Satisfies Differential Coulomb Field Law]
  \label{def:physics:phyx-mini-0995:phyxminiproblems-problemphyxmini0995-satisfiesdifferentialcoulombfieldlaw}
  \lean{PhyXMiniProblems.ProblemPhyXMini0995.SatisfiesDifferentialCoulombFieldLaw}
  \uses{def:physics:phyx-mini-0995:phyxminiproblems-problemphyxmini0995-angularchargedensityincoulombsperradian, def:physics:phyx-mini-0995:phyxminiproblems-problemphyxmini0995-chargedringaxissetup, def:physics:phyx-mini-0995:phyxminiproblems-problemphyxmini0995-sourcetoobservationdisplacementinmeters, def:physics:phyx-mini-0995:phyxminiproblems-problemphyxmini0995-sourcetoobservationdistanceinmeters}
  Differential Coulomb's law in coherent SI units: dE/dφ = k (dQ/dφ) r⃗ / ‖r⃗‖³. This general inverse-square vector law retains both axial and transverse contributions and is not the closed-form field requested by the question
\end{definition}
\begin{proof}
  This is the declared model definition of Satisfies Differential Coulomb Field Law.
\end{proof}
\begin{definition}[Satisfies Ring Field Superposition]
  \label{def:physics:phyx-mini-0995:phyxminiproblems-problemphyxmini0995-satisfiesringfieldsuperposition}
  \lean{PhyXMiniProblems.ProblemPhyXMini0995.SatisfiesRingFieldSuperposition}
  \uses{def:physics:phyx-mini-0995:phyxminiproblems-problemphyxmini0995-chargedringaxissetup}
  The net electric field at P is the angular integral of all dE vectors
\end{definition}
\begin{proof}
  This is the declared model definition of Satisfies Ring Field Superposition.
\end{proof}
\begin{lemma}[source To Observation Distance eq]
  \label{lem:physics:phyx-mini-0995:phyxminiproblems-problemphyxmini0995-sourcetoobservationdistance-eq}
  \lean{PhyXMiniProblems.ProblemPhyXMini0995.sourceToObservationDistance_eq}
  \uses{def:physics:phyx-mini-0995:phyxminiproblems-problemphyxmini0995-lengthinmeters, def:physics:phyx-mini-0995:phyxminiproblems-problemphyxmini0995-chargedringaxissetup, def:physics:phyx-mini-0995:phyxminiproblems-problemphyxmini0995-sourcetoobservationdistanceinmeters, def:physics:phyx-mini-0995:phyxminiproblems-problemphyxmini0995-haschargedringaxisgeometry}
  Every point on the parametrized ring is the displayed distance √(x² + a²) from P
\end{lemma}
\begin{proof}
  The conclusion follows from the governing relations and figure readouts named in the statement of source To Observation Distance eq.
\end{proof}
\begin{lemma}[transverse Field Components cancel]
  \label{lem:physics:phyx-mini-0995:phyxminiproblems-problemphyxmini0995-transversefieldcomponents-cancel}
  \lean{PhyXMiniProblems.ProblemPhyXMini0995.transverseFieldComponents_cancel}
  \uses{def:physics:phyx-mini-0995:phyxminiproblems-problemphyxmini0995-chargedringaxissetup, def:physics:phyx-mini-0995:phyxminiproblems-problemphyxmini0995-haschargedringaxisgeometry, def:physics:phyx-mini-0995:phyxminiproblems-problemphyxmini0995-satisfiesuniformringdistribution, def:physics:phyx-mini-0995:phyxminiproblems-problemphyxmini0995-satisfiesdifferentialcoulombfieldlaw, def:physics:phyx-mini-0995:phyxminiproblems-problemphyxmini0995-satisfiesringfieldsuperposition}
  Opposite source elements have equal-and-opposite y and z field components, so the complete ring field has no transverse component
\end{lemma}
\begin{proof}
  The conclusion follows from the governing relations and figure readouts named in the statement of transverse Field Components cancel.
\end{proof}
\begin{lemma}[axial Field Component eq]
  \label{lem:physics:phyx-mini-0995:phyxminiproblems-problemphyxmini0995-axialfieldcomponent-eq}
  \lean{PhyXMiniProblems.ProblemPhyXMini0995.axialFieldComponent_eq}
  \uses{def:physics:phyx-mini-0995:phyxminiproblems-problemphyxmini0995-chargeincoulombs, def:physics:phyx-mini-0995:phyxminiproblems-problemphyxmini0995-lengthinmeters, def:physics:phyx-mini-0995:phyxminiproblems-problemphyxmini0995-chargedringaxissetup, def:physics:phyx-mini-0995:phyxminiproblems-problemphyxmini0995-haschargedringaxisgeometry, def:physics:phyx-mini-0995:phyxminiproblems-problemphyxmini0995-hasphysicalelectrostaticparameters, def:physics:phyx-mini-0995:phyxminiproblems-problemphyxmini0995-satisfiesuniformringdistribution, def:physics:phyx-mini-0995:phyxminiproblems-problemphyxmini0995-satisfiesdifferentialcoulombfieldlaw, def:physics:phyx-mini-0995:phyxminiproblems-problemphyxmini0995-satisfiesringfieldsuperposition}
  All axial contributions have the same component. Integrating the uniform charge density gives the standard scalar coefficient of the ring field
\end{lemma}
\begin{proof}
  The conclusion follows from the governing relations and figure readouts named in the statement of axial Field Component eq.
\end{proof}
\begin{definition}[Answer Choice]
  \label{def:physics:phyx-mini-0995:phyxminiproblems-problemphyxmini0995-answerchoice}
  \lean{PhyXMiniProblems.ProblemPhyXMini0995.AnswerChoice}
  Labels attached to the four electric-field choices in the source
\end{definition}
\begin{proof}
  This is the declared model definition of Answer Choice.
\end{proof}
\begin{definition}[axial Coefficient]
  \label{def:physics:phyx-mini-0995:phyxminiproblems-problemphyxmini0995-answerchoice-axialcoefficient}
  \lean{PhyXMiniProblems.ProblemPhyXMini0995.AnswerChoice.axialCoefficient}
  \uses{def:physics:phyx-mini-0995:phyxminiproblems-problemphyxmini0995-chargeincoulombs, def:physics:phyx-mini-0995:phyxminiproblems-problemphyxmini0995-lengthinmeters, def:physics:phyx-mini-0995:phyxminiproblems-problemphyxmini0995-chargedringaxissetup, def:physics:phyx-mini-0995:phyxminiproblems-problemphyxmini0995-answerchoice}
  The coherent-SI scalar multiplying î in each printed answer. This records the alternatives but does not define the independent electric field
\end{definition}
\begin{proof}
  This is the declared model definition of axial Coefficient.
\end{proof}
\begin{definition}[field Vector]
  \label{def:physics:phyx-mini-0995:phyxminiproblems-problemphyxmini0995-answerchoice-fieldvector}
  \lean{PhyXMiniProblems.ProblemPhyXMini0995.AnswerChoice.fieldVector}
  \uses{def:physics:phyx-mini-0995:phyxminiproblems-problemphyxmini0995-chargedringaxissetup, def:physics:phyx-mini-0995:phyxminiproblems-problemphyxmini0995-xaxisunit, def:physics:phyx-mini-0995:phyxminiproblems-problemphyxmini0995-answerchoice}
  The vector represented by a displayed answer choice
\end{definition}
\begin{proof}
  This is the declared model definition of field Vector.
\end{proof}
\begin{definition}[recorded Dataset Answer]
  \label{def:physics:phyx-mini-0995:phyxminiproblems-problemphyxmini0995-recordeddatasetanswer}
  \lean{PhyXMiniProblems.ProblemPhyXMini0995.recordedDatasetAnswer}
  \uses{def:physics:phyx-mini-0995:phyxminiproblems-problemphyxmini0995-answerchoice}
  The dataset's recorded answer label, retained as unused metadata
\end{definition}
\begin{proof}
  This is the declared model definition of recorded Dataset Answer.
\end{proof}
% --- Archon declaration topology end ---
% --- Archon physics formalization source end ---
